% Supplementary Material: proofs
\documentclass[11pt,a4paper]{article}
\usepackage[utf8]{inputenc}
\usepackage[T1]{fontenc}
\usepackage{lmodern}
\usepackage[a4paper,margin=2.5cm]{geometry}
\usepackage{amsmath,amssymb,amsthm,mathtools}
\usepackage{bm}
\usepackage{booktabs}
\usepackage{enumitem}
\usepackage{xcolor}
\usepackage{listings}
\usepackage[natbibapa]{apacite}
\usepackage{aliascnt}

% ---- listing style ----
\definecolor{kw}{RGB}{30,64,175}
\definecolor{cmt}{RGB}{22,101,52}
\definecolor{strc}{RGB}{159,18,57}
\definecolor{lstbg}{RGB}{250,250,250}
\lstdefinelanguage{cfc}{
  morekeywords={floor,circuit,species,reaction,solve,yield,patient,probe,
    panel,measure,localize,using,until,resolved,otherwise,decline,cell,
    gap,close,toward,coherent,basis,assert,emit,report,let,import,
    holonomy,loops,of,wegscheider,check,from,where,in,foreach,scan,
    against,collect,tolerance,witness,admit,invalid,negative},
  morekeywords=[2]{Reading,Cell,Gap,Probe,Panel,Verdict,Concern,Circuit,
    Loop,Holonomy,HolonomyVector,Tolerance,Residual},
  sensitive=true,
  morecomment=[l]{--},
  morestring=[b]",
  keywordstyle=\color{kw}\bfseries,
  keywordstyle=[2]\color{strc},
  commentstyle=\color{cmt}\itshape,
  stringstyle=\color{strc},
  basicstyle=\ttfamily\small,
  backgroundcolor=\color{lstbg},
  frame=single,framesep=4pt,xleftmargin=6pt,xrightmargin=4pt,
  breaklines=true,showstringspaces=false,columns=fullflexible,tabsize=2
}
\lstset{language=cfc}

\usepackage{hyperref}
\usepackage{cleveref}
\usepackage{xr}

% Import the main manuscript's labels so that every \Cref in a proof
% resolves to the number the reader sees in the main text.
\externaldocument{medinformatics-submission}

% Theorem environments, so \Cref names ("Theorem", "Lemma", ...) resolve.
% Aliased counters: shared numbering, but cleveref sees each type.
\newcommand{\newsharedtheorem}[2]{%
  \newaliascnt{#1}{theorem}\newtheorem{#1}[#1]{#2}\aliascntresetthe{#1}}
\newtheorem{theorem}{Theorem}[section]
% Sections here are unnumbered, so number results S1, S2, ... to keep
% them distinct from the main text's numbering. Set before the aliases,
% which copy this format when they are created.
\renewcommand{\thetheorem}{S\arabic{theorem}}
\renewcommand{\thetable}{S\arabic{table}}
\newsharedtheorem{lemma}{Lemma}
\newsharedtheorem{proposition}{Proposition}
\newsharedtheorem{corollary}{Corollary}
\theoremstyle{definition}
\newsharedtheorem{definition}{Definition}
\newsharedtheorem{assumption}{Assumption}
\newsharedtheorem{criterion}{Criterion}
\newsharedtheorem{protocol}{Protocol}
\newsharedtheorem{rulename}{Rule}
\theoremstyle{remark}
\newsharedtheorem{remark}{Remark}
\crefname{assumption}{Assumption}{Assumptions}
\Crefname{assumption}{Assumption}{Assumptions}
\crefname{criterion}{Criterion}{Criteria}
\Crefname{criterion}{Criterion}{Criteria}
\crefname{protocol}{Protocol}{Protocols}
\Crefname{protocol}{Protocol}{Protocols}
\crefname{rulename}{Rule}{Rules}
\Crefname{rulename}{Rule}{Rules}

\newcommand{\R}{\mathbb{R}}
\newcommand{\abs}[1]{\lvert#1\rvert}
\newcommand{\norm}[1]{\lVert#1\rVert}
\newcommand{\eps}{\varepsilon}
\newcommand{\epsnum}{\eps_{\mathrm{num}}}
\newcommand{\epsdata}{\eps_{\mathrm{data}}}
\newcommand{\epsstar}{\eps^{\ast}}
\newcommand{\hol}{H}
\newcommand{\mach}{u}
\newcommand{\cf}[1]{\texttt{#1}}

\title{Supplementary Material\\[4pt]
\large Proofs for ``Separating Numerical Noise from Thermodynamic
Inconsistency in Cycle-Consistency Tests on Metabolic Circuits''}
\author{Kundai Farai Sachikonye}
\date{}

\begin{document}
\maketitle

\noindent
This document collects the proofs of every numbered result in the main
text, in order of appearance. Section~S$k$ here proves the result whose
proof the main text defers to Section~S$k$. Notation, numbering, and
equation references follow the main text.

\section*{S1}
\addcontentsline{toc}{section}{S1}
(i) The net rate of change of $c_i$ is incoming minus outgoing
conversion rate; under \eqref{eq:GJ} the conversion rate along $(i,j)$
is $J_{ij}$, so stationarity of $c_i$ is exactly the vanishing of the
signed sum, which is the current law at node $i$.

(ii) The sum telescopes to $\mu_{i_1}-\mu_{i_1}=0$ identically, because
$\mu$ is a function on vertices. The substantive content is the
converse: a set of rate constants admits a consistent potential function
$\mu$ if and only if the products of forward rate constants around every
cycle equal the products of reverse constants --- the Wegscheider
condition \citep{wegscheider1901,schuster1989}. A network violating it
admits no $\mu$ and hence no circuit representation.

\section*{S2}
\addcontentsline{toc}{section}{S2}
$z_\ell\in\ker B$ by \Cref{def:cyclespace}, so
$B z_\ell = 0$ and hence $\hol_\ell(\mu)=z_\ell^{\top}B^{\top}\mu
= (Bz_\ell)^{\top}\mu = 0$.

\section*{S3}
\addcontentsline{toc}{section}{S3}
$\operatorname{im}B^{\top} = (\ker B)^{\perp} = \mathcal{Z}(G)^{\perp}$
by the fundamental theorem of linear algebra. Hence $d\in
\operatorname{im}B^{\top}$ iff $z^{\top}d=0$ for all
$z\in\mathcal{Z}(G)$, and since a basis spans $\mathcal{Z}(G)$, iff
$z_\ell^{\top}d=0$ for every $\ell$ in that basis.

\section*{S4}
\addcontentsline{toc}{section}{S4}
Each difference $\Delta\mu_e = \mu_j-\mu_i$ is computed with one
rounding: $\mathrm{fl}(\mu_j-\mu_i) = (\mu_j-\mu_i)(1+\theta_e)$,
$\abs{\theta_e}\le\mach$, giving an absolute error at most
$\mach\abs{\mu_j-\mu_i}\le 2\mach\Lambda_\ell$. Summing $L$ such
quantities sequentially, the standard forward error bound for
floating-point summation \citep[Ch.~4]{higham2002} gives a total error
at most $\gamma_{L}\sum_{e\in\ell}\abs{\Delta\mu_e}$ from the additions,
plus the propagated per-difference errors, at most
$\sum_{e}\mach\abs{\Delta\mu_e}$. Since
$\abs{\Delta\mu_e}\le 2\Lambda_\ell$ and there are $L$ edges,
$\sum_{e\in\ell}\abs{\Delta\mu_e}\le 2L\Lambda_\ell$. Collecting the two
contributions and absorbing the constants into $\gamma_{2L}$ yields
\eqref{eq:noise-bound}. Because the exact value is $0$, the computed
value equals its own error, which is what the display asserts.

For the compensated variant, Kahan summation \citep{kahan1965,
higham2002} replaces the $\gamma_L$ factor by $2\mach + O(L\mach^{2})$,
leaving only the $L$ per-difference roundings, hence
$O(L\mach\Lambda_\ell)$.

The bound is stated in $\Lambda_\ell=\max\abs{\mu_i}$ only through
$\abs{\Delta\mu_e}\le 2\Lambda_\ell$; every step above holds with
$2\Lambda_\ell$ replaced by the spread $\max\mu-\min\mu$ on the cycle,
which is gauge-invariant. When all potentials on the cycle lie within a
factor of two of one another, Sterbenz's lemma \citep{goldberg1991} makes
each difference exact, and since each partial sum
$\mu_{i_{k+1}}-\mu_{i_1}$ is then itself a difference of two such
potentials, every addition is exact as well: the computed cycle sum is
exactly zero.

\section*{S5}
\addcontentsline{toc}{section}{S5}
Immediate from \eqref{eq:noise-bound}, whose right-hand side is, for
fixed $\mach$, proportional to $\gamma_{2L}\,L\Lambda_\ell$ and hence to
$L^{2}\Lambda_\ell$ to first order in $L\mach$. Two cycles whose values of
$L^{2}\Lambda$ differ by a factor $r$ have floors differing by $r$, and
a constant threshold separating noise from signal on both must lie
within a factor $r$ of each floor at once, which fails once $r$ exceeds
the dynamic range the test requires.

\section*{S6}
\addcontentsline{toc}{section}{S6}
$\hol_\ell(\mu+c\mathbf{1}) = z_\ell^{\top}B^{\top}(\mu+c\mathbf{1})
= \hol_\ell(\mu) + c\,z_\ell^{\top}B^{\top}\mathbf{1} = \hol_\ell(\mu)$,
since $B^{\top}\mathbf{1}=0$ (every edge has one head and one tail).
The stated $c^{\ast}$ centres the interval $[\min\mu,\max\mu]$ on zero,
minimising $\max\abs{\mu_i}$, which is the definition of $\Lambda_\ell$.

\section*{S7}
\addcontentsline{toc}{section}{S7}
\emph{\Cref{lem:species-cancel}.} With $d=B^{\top}\hat\mu$,
$\hol_\ell(d)=z_\ell^{\top}B^{\top}\hat\mu=(Bz_\ell)^{\top}\hat\mu=0$,
since $Bz_\ell=0$ by \Cref{def:cyclespace}. The argument holds for every
$\hat\mu$, and in particular for $\hat\mu=\mu+\eta$ with any $\eta$.
Equivalently, the coefficient of $\hat\mu_i$ in the cycle sum is the
number of times the cycle enters species $i$ minus the number of times
it leaves, which is zero for every species of a closed walk.

\emph{\Cref{thm:data-floor}.} With $d=\bar d+\eta$ and $\bar d$
consistent, $\hol_\ell(d)=z_\ell^{\top}\bar d+z_\ell^{\top}\eta
=z_\ell^{\top}\eta$ by \Cref{prop:detect-defect}. This is linear in the
independent errors, so it has mean zero and variance
$\sum_e z_\ell(e)^{2}\sigma_e^{2}$, which is \eqref{eq:data-sd}; for a
simple cycle $z_\ell(e)\in\{-1,0,1\}$ and the sum runs over the edges of
$\ell$. Under normality $z_\ell^{\top}\eta$ is normal, and
$\Pr[\abs{z_\ell^{\top}\eta}\le z_{1-\alpha/2}s_\ell]=1-\alpha$ exactly.
Without normality, Chebyshev's inequality gives the same statement with
$z_{1-\alpha/2}$ replaced by $\alpha^{-1/2}$.

The two results are the two halves of one fact: the cycle sum reads only
the part of the edge data that is not a gradient, so an error enters it
if and only if it enters edge by edge.

\section*{S8}
\addcontentsline{toc}{section}{S8}
With exact data, \Cref{lem:vanish} gives $\hol_\ell = 0$ and
\Cref{thm:noise-scale} bounds the computed value by $\epsnum(\ell) \le
\epsstar(\ell)$; the strict inequality in the event is therefore never
satisfied. With uncertain edge data, $\abs{\hol_\ell} \le \epsdata(\ell)$
with probability at least $1-\alpha$ by \Cref{thm:data-floor}, and on
that event the computed value satisfies
$\abs{\widehat{\hol}_\ell}\le\abs{\hol_\ell}+\epsnum(\ell)
\le\epsdata(\ell)+\epsnum(\ell)=\epsstar(\ell)$. Hence
$\Pr[\abs{\widehat{\hol}_\ell}>\epsstar(\ell)]\le\alpha$. (With
$\epsstar$ the maximum of the two floors rather than their sum, the last
inequality fails, which is why the sum is used.) Strictly, the rounding
bound is here applied to a sum whose exact value is $\hol_\ell$ rather
than zero; the argument of S4 applies unchanged with
$\sum_{e\in\ell}\abs{d_e}$ in place of $2L\Lambda_\ell$, which for edge
data within a few $\sigma_e$ of a gradient changes the bound
negligibly.

\section*{S9}
\addcontentsline{toc}{section}{S9}
With exact data, the computed value differs from the true one by at most
$\epsnum(\ell)\le\epsstar(\ell)$ (\Cref{thm:noise-scale}), so
$\abs{\widehat{\hol}_\ell} \ge D - \epsstar(\ell) > 2\epsstar(\ell) -
\epsstar(\ell) = \epsstar(\ell)$. With independent errors on the edge
data, the true sum is $\pm D+z_\ell^{\top}\eta$, and
$\abs{z_\ell^{\top}\eta}\le\epsdata(\ell)$ with probability at least
$1-\alpha$ (S7). On that event
$\abs{\widehat{\hol}_\ell}\ge D-\epsdata(\ell)-\epsnum(\ell)
=D-\epsstar(\ell)>\epsstar(\ell)$.

\section*{S10}
\addcontentsline{toc}{section}{S10}
(i) The three conditions partition $[0,\infty)$ at the two cut points
$\epsnum \le \epsstar = \epsnum+\epsdata$. When $\epsdata$ is undefined
(\Cref{rem:no-sigma}) there is one cut point and every value above it is
\textsf{UNDECIDABLE}.

(ii) The middle interval $(\epsnum,\epsnum+\epsdata]$ is non-empty
exactly when $\epsdata>0$. Numerically, the measured noise in a long
uncentred high-potential cycle is of order $10^{-12}$ kJ/mol and the
bound $\epsnum$ at most of order $10^{-9}$ (main text, \Cref{tab:v1} and
\Cref{fig:tol_panel1}), while edge-level uncertainties of $1$--$10$
kJ/mol \citep{noor2013,flamholz2012} give $\epsdata$ of order $1$--$10$
kJ/mol. The middle interval therefore spans some nine to thirteen orders
of magnitude.

(iii) A two-valued rule places a single cut at some $t$. If
$t<\epsstar$, cycles with $\abs{\widehat{\hol}}\in(t,\epsstar]$ are
called inconsistent though their sums are within what the data's own
uncertainty and the arithmetic can produce from a consistent network. If
$t\ge\epsstar$, then $t>\epsnum$, and cycles with
$\abs{\widehat{\hol}}\in(\epsnum,t]$ are called consistent though their
sums exceed what the arithmetic alone can produce from exact data, so
the verdict asserts a consistency the data do not establish. Either way
one of the two misstatements occurs.

\section*{S11}
\addcontentsline{toc}{section}{S11}
Each $\mathcal{B}_i$ spans $\mathcal{Z}(G)$ by \Cref{def:cyclespace}.
The map $z\mapsto z^{\top}d$ is linear, so it vanishes on a spanning set
if and only if it vanishes on the whole space, and hence if and only if
it vanishes on any other spanning set.

\section*{S12}
\addcontentsline{toc}{section}{S12}
Let $G$ be the ``theta-plus'' graph: vertices $\{a,b\}$ joined by four
internally disjoint paths $P_1,\dots,P_4$, giving
$\nu = 3$. Let $e^{\ast}\in P_1$ carry a defect of size $D$, all other
edges consistent, so that any cycle containing $e^{\ast}$ exactly once
has cycle sum $\pm D$ and any cycle avoiding it has sum $0$.

Take $\mathcal{B}_1 = \{P_1P_2^{-1},\,P_2P_3^{-1},\,P_3P_4^{-1}\}$.
Only the first contains $e^{\ast}$, so exactly one element has nonzero
sum.

Take $\mathcal{B}_2 = \{P_1P_2^{-1},\,P_1P_3^{-1},\,P_1P_4^{-1}\}$.
This is a basis: the three are independent, since $P_1P_j^{-1}$ is the
only element containing $P_j$, for $j=2,3,4$. Every element
contains $e^{\ast}$, so all three have nonzero sum. Both are cycle bases
of the same graph and the same data, and they flag one and three cycles
respectively.

\section*{S13}
\addcontentsline{toc}{section}{S13}
With the remaining edge data exact, a cycle's true sum is $D$ times its
net signed traversal count of $e^{\ast}$, and $0$ from every other edge.
A cycle with nonzero count has $\abs{\hol_\ell}\ge D>2\epsstar(\ell)$
and is flagged by \Cref{thm:detection}; a cycle with zero count has
$\hol_\ell=0$ and is not flagged, by \Cref{thm:no-false-positive}. For
simple cycles a nonzero count is containment. Hence every flagged cycle
contains $e^{\ast}$, so $e^{\ast}$ lies in their intersection, which is
$W(\mathcal{B})$. With independent errors on the edge data, each of the
$\nu$ basis cycles receives the verdict the exact case assigns it except
with probability at most $\alpha$ (\Cref{thm:no-false-positive,%
thm:detection}); by the union bound all $\nu$ do so with probability at
least $1-\nu\alpha$, and on that event the argument above applies.

\section*{S14}
\addcontentsline{toc}{section}{S14}
By \Cref{rule:verdict} the only introduction form for \cf{Verdict}
requires a premise of type \cf{Tolerance}, and the conclusion carries it
by \Cref{def:types}. A numeric literal has no type in the grammar of
\Cref{def:types}, so there is no derivation in which one appears in
the second premise; and no coercion rule from a numeric literal to
\cf{Tolerance} exists, since \cf{Tolerance} is introduced only by the
computation of \eqref{eq:epsstar} from a \cf{Loop} and its data. Hence
no derivation produces a bare verdict.


%% ---- moved from the main text for length ----

\section*{S15: a language in which the tolerance cannot be omitted}
\addcontentsline{toc}{section}{S15}
% =====================================================================

The analysis of the main text is only as good as its
implementation. A practitioner who computes $\epsstar$ correctly and then writes
\verb|if abs(H) > 1e-6| has gained nothing. We therefore express the
protocol in a small typed language in which the omission is not
expressible.

\subsection*{Types}

\begin{definition}[Types]
\label{def:types}
The types are
\[
  \tau \;::=\; \cf{Circuit} \mid \cf{Loop} \mid \cf{Holonomy}
  \mid \cf{Tolerance} \mid \cf{Verdict} \mid \cf{Witness}.
\]
A value of type \cf{Holonomy} is a pair $(\widehat{\hol}_\ell,\ell)$; a
value of type \cf{Tolerance} is a triple
$(\epsnum,\epsdata,\epsstar)$ with $\epsdata$ possibly the distinguished
value \cf{undefined}; a value of type \cf{Verdict} is an element of
$\{\textsf{CONSISTENT},\textsf{UNDECIDABLE},\textsf{INCONSISTENT}\}$
\emph{paired with} the \cf{Tolerance} that produced it.
\end{definition}

\subsection*{The typing rule that does the work}

\begin{rulename}[Verdict requires a tolerance]
\label{rule:verdict}
\[
  \frac{\Gamma\vdash h : \cf{Holonomy}
        \qquad \Gamma\vdash t : \cf{Tolerance}}
       {\Gamma\vdash \cf{admit}\ h\ \cf{tolerance}\ t : \cf{Verdict}}
\]
There is no other rule with \cf{Verdict} in its conclusion.
\end{rulename}

\begin{theorem}[No bare verdict]
\label{thm:no-bare-verdict}
No well-typed program produces a value of type \cf{Verdict} that is not
paired with the \cf{Tolerance} used to derive it. In particular, a
comparison of a \cf{Holonomy} against a numeric literal is not
well-typed.
\end{theorem}

\noindent\emph{Proof.} See Section~S14 above.

\begin{remark}[What this does and does not buy]
\label{rem:type-limits}
\Cref{thm:no-bare-verdict} rules out a class of implementation error:
it is not possible to forget the tolerance, to hard-code it, or to
report a verdict whose provenance is not recoverable. It does not make
the tolerance \emph{correct}: a program supplying a wrong $\sigma_i$ to
\eqref{eq:data-floor} is well-typed and wrong. The type system enforces
accountability of the computation, not accuracy of the inputs, and we
state this because the distinction is easy to overstate in the other
direction.
\end{remark}

\subsection*{The protocol as a program}

\begin{lstlisting}[caption={The cycle-consistency test with cycle-local
tolerances. Every verdict carries its tolerance by
\Cref{thm:no-bare-verdict}; the validity gate of \Cref{crit:invalid}
precedes any diagnostic claim.},label={lst:program}]
-- kcl-tolerance.cfc
-- Cycle-consistency with per-cycle tolerances.

import Recon3D
import Recon3D.uncertainty as sigma

circuit Reference from Recon3D.glycolysis_tca {
  solve yield C_ref
}

-- Gauge-centre potentials to tighten the bound on the numerical
-- floor (gauge proposition). Changes no cycle sum.
let C_ref_centred := centre_potentials(C_ref)

-- Minimum-weight cycle basis: sharper witness sets.
let B := minimum_cycle_basis(C_ref_centred)

-- VALIDITY GATE (invalid versus negative). The reference network
-- must pass before any diagnostic verdict is meaningful. Data floors
-- here are taken at level alpha / size(B).
foreach loop in B {
  holonomy of loop in C_ref_centred yield h_ref
  tolerance of loop with sigma yield t
  admit h_ref tolerance t yield v_ref

  assert v_ref != INCONSISTENT
    otherwise invalid
      emit "reference network fails its own consistency check"
}

emit "validity gate passed: reference is self-consistent"

-- DIAGNOSTIC. Perturbed network: the forward constant alone is
-- scaled, which breaks the Haldane relation and injects a defect of
-- RT ln(1/0.15) = 4.9 kJ/mol on the PK edge.
circuit Perturbed from Recon3D.glycolysis_tca {
  reaction PK : PEP -> Pyruvate, k : 0.15 * Recon3D.k_PK
  solve yield C_pert
}

let C_pert_centred := centre_potentials(C_pert)

let flagged := {}
foreach loop in B {
  holonomy of loop in C_pert_centred yield h
  tolerance of loop with sigma yield t
  admit h tolerance t yield v

  report loop, h, t.numerical, t.data, t.star, v

  where v == INCONSISTENT collect loop into flagged
}

-- Localisation is to the witness set, not to a named loop.
witness of flagged yield W
report "witness_set", W
report "witness_cardinality", size(W)

-- The claim under test, stated basis-independently.
assert Recon3D.edge("PK") in W
  emit "defect localised: PK edge in witness set"
  otherwise decline
    emit "defect not localised at this data quality"
\end{lstlisting}

% =====================================================================
\section*{S16: protocol and pre-registered criteria}
\addcontentsline{toc}{section}{S16}
% =====================================================================

\subsection*{Data}

\begin{enumerate}[leftmargin=1.5em,itemsep=2pt]
\item \textbf{Network.} The glycolysis, pentose phosphate, and TCA
subnetwork of a human metabolic reconstruction
\citep{brunk2018}, approximately $80$ reactions over $60$ metabolites.
\item \textbf{Edge data.} For each reaction with forward and reverse rate
constants in the kinetic models \citep{rapoport1975,mulquiney1999}, the
edge datum is the standard reaction energy those constants imply by the
Haldane relation, $d_e=-RT\ln(k_e^{+}/k_e^{-})$, with its reported
uncertainty $\sigma_e$; where an equilibrium constant is reported
independently for the reaction, it is used instead. Edges whose datum is
derived from component-contribution formation energies
\citep{noor2013,flamholz2012} are consistent by construction
(\Cref{lem:species-cancel}) and are recorded as such. Edges without a
reported uncertainty are recorded and trigger \Cref{rem:no-sigma}.
\item \textbf{Reference physiology.} Erythrocyte metabolite
concentrations and glycolytic flux from established kinetic models
\citep{rapoport1975,mulquiney1999}.
\item \textbf{Perturbation.} The pyruvate kinase forward rate constant
reduced to $15\%$ of reference with the reverse constant unchanged,
which injects a defect of $RT\ln(1/0.15)\approx4.9$ kJ/mol on the PK
edge. The $15\%$ matches the residual activity documented in the
clinical deficiency \citep{zanella2005}; the deficiency itself scales
both directions, which is Control~1.
\end{enumerate}

\subsection*{Procedure}

\begin{protocol}[Cycle-consistency with cycle-local tolerance]
\leavevmode
\begin{enumerate}[leftmargin=1.5em,itemsep=3pt]
\item Construct the reference circuit by
\Cref{def:mu,def:conductance}; verify the current law at every internal
node to the numerical floor of \Cref{thm:noise-scale} applied to the
node balance.
\item Gauge-centre the potentials (\Cref{prop:gauge}) and record the
resulting $\Lambda_\ell$ per cycle, both before and after.
\item Compute a minimum-weight cycle basis \citep{horton1987}; record
$\nu$, the length distribution, and the basis itself so that the
computation is reproducible. Verify that the PK edge lies on at least
one cycle; if it is a bridge, no cycle sum can detect the perturbation
and the experiment is invalid by design.
\item For each cycle compute $\epsnum(\ell)$ by \eqref{eq:noise-bound}
and $\epsdata(\ell)$ by \eqref{eq:data-floor}, at $\alpha/\nu$ for the
validity gate and at $\alpha=0.01$ otherwise; record both, and
$\epsstar$.
\item \textbf{Validity gate.} Evaluate \Cref{def:validity} on the
reference. If it fails, halt and report an \emph{invalid} experiment per
\Cref{crit:invalid}.
\item Evaluate the perturbed circuit; record $\widehat{\hol}_\ell$ and
$\mathrm{verdict}(\ell)$ for every cycle.
\item Compute the witness set $W$ (\Cref{def:witness}) and record
$\abs{W}$ and whether the perturbed edge lies in it.
\item \textbf{Control 1 (activity loss).} Repeat step 6 with both PK
constants scaled by $0.15$, which preserves the Haldane relation and
every cycle sum. Any \textsf{INCONSISTENT} verdict here is a false
positive; at level $\alpha/\nu$ per cycle its probability is at most
$\alpha$ (\Cref{thm:no-false-positive}). This control also shows
directly that the test is blind to a loss of activity that leaves the
equilibrium constant unchanged.
\item \textbf{Control 2 (basis).} Repeat steps 3--7 with a
\emph{different} cycle basis (e.g.\ a fundamental basis from a spanning
tree). Verify the network-level verdict agrees
(\Cref{thm:basis-verdict}) and record how $F$ and $W$ differ
(\Cref{prop:basis-localisation}).
\item \textbf{Control 3 (detection sweep).} Sweep the perturbation
magnitude and record the smallest $D$ at which the correct edge enters
$W$; compare against the predicted $2\epsstar$ of
\Cref{thm:detection}.
\end{enumerate}
\end{protocol}

\subsection*{Pre-registered criteria}

\begin{criterion}[Success]
\label{crit:success}
\leavevmode
\begin{enumerate}[label=(\alph*),leftmargin=2em,itemsep=2pt]
\item \textbf{Validity.} The reference network passes
\Cref{def:validity} on every basis cycle.
\item \textbf{No false positives.} Control 1 produces zero
\textsf{INCONSISTENT} verdicts, as \Cref{thm:no-false-positive}
requires.
\item \textbf{Localisation.} The perturbed edge lies in $W$, and
$\abs{W}\le 3$.
\item \textbf{Basis invariance.} Control 2 yields the same network-level
verdict under both bases.
\item \textbf{Detection threshold.} Control 3 finds the empirical
detection threshold within a factor of $10$ of the predicted
$2\epsstar$.
\end{enumerate}
\end{criterion}

\begin{criterion}[Failure, and what each failure means]
\label{crit:failure}
\leavevmode
\begin{itemize}[leftmargin=1.5em,itemsep=2pt]
\item Failure of (a) renders the experiment \emph{invalid}
(\Cref{crit:invalid}): the reconstruction's thermodynamic annotations
are not self-consistent to their claimed precision. No conclusion about
the diagnostic follows.
\item Failure of (b) has probability at most $\alpha$ under the stated
error model. It indicates an error in the tolerance computation, or
reported uncertainties that understate the real errors of the edge
data; it is not a detection.
\item Failure of (c) with (a) and (b) passing is a genuine
\emph{negative}: the test does not localise this defect at this data
quality. If additionally most verdicts are \textsf{UNDECIDABLE}, the
limiting factor is $\epsdata$ and the finding is about the database
(\Cref{rem:undecidable}).
\item Failure of (d) would contradict \Cref{thm:basis-verdict} and
indicates an implementation error in the basis computation.
\item Failure of (e) indicates the bound of \Cref{thm:noise-scale} is
loose by more than an order of magnitude on real data, which is a
finding about the bound worth reporting.
\end{itemize}
\end{criterion}

\begin{remark}[Controls that can fail]
\label{rem:controls}
Controls 1--3 are chosen so that each can fail on well-formed input.
Control 1 tests a theorem whose failure probability is bounded by
$\alpha$, so a failure is strong evidence against the computation or the
error model. Control 2 tests an invariance that the theory asserts, so
a failure localises to the implementation. Control 3 tests a bound
rather than an identity, so a failure is informative about the bound's
tightness. A control that cannot fail measures nothing, and we have
excluded from the protocol several checks with that property --- notably
any check comparing the implementation's verdict against the same
implementation's tolerance, which would pass by construction.
\end{remark}


\subsection*{Outcome on the erythrocyte network}

We executed the protocol on the human erythrocyte network assembled from
UniProtKB/Swiss-Prot \citep{uniprot2025} and Rhea \citep{bansal2022}
(main text, \Cref{sec:data}) rather than on the Recon3D glycolysis--TCA
subnetwork named under \emph{Data} above, since erythrocytes have no
mitochondria and no tricarboxylic acid cycle. Step~3 fails: the pyruvate
kinase reaction lies on none of the network's seven cycles, because no
erythrocyte enzyme returns pyruvate to phosphoenolpyruvate. By the
criterion of step~3 the experiment is invalid by design, and steps
4--10 were not run. No erythrocyte cycle has a measured equilibrium
constant on every edge (Section~S17), so on present data there is no
substitute target either.

\section*{S17: analysis of experimental data}

The code is \texttt{validation/realdata/run\_realdata.py}; every number
below is read from \texttt{results\_real.json}.

\subsection*{Cycles of a stoichiometric network}

Let $S\in\mathbb{Z}^{n\times m}$ be the stoichiometric matrix of $m$
reactions over $n$ species, after balancing with water and protons and
removing the proton row. A cycle is an integer $z\in\ker S$; its sum on
edge data $d\in\mathbb{R}^{m}$ is $\hol_z=z^{\top}d$. If
$d=S^{\top}\mu$ for any $\mu\in\mathbb{R}^{n}$, then
$\hol_z=(Sz)^{\top}\mu=0$; this is \Cref{lem:vanish}, and with $\mu$
replaced by species-level errors it is \Cref{lem:species-cancel}. The
proof of \Cref{thm:data-floor} uses only the linearity of $\hol_z$ in
$d$ and so holds unchanged, with variance $\sum_e z_e^{2}\sigma_e^{2}$.
The incidence matrix of the main text is the case in which every column
of $S$ has one $+1$ and one $-1$.

\begin{proposition}[Numerical floor for edge data]
\label{prop:edge-floor}
Let the edge data $d$ be held in the arithmetic of \Cref{ass:fp}, and
let $z$ have $k$ nonzero entries, each exactly representable. Then the
value $\widehat{\hol}_z$ of $z^{\top}d$ computed by forming the $k$
products and summing them in any order satisfies
\[
  \bigl\lvert\widehat{\hol}_z-\hol_z\bigr\rvert
  \;\le\;\gamma_k\sum_{e}\abs{z_e d_e}.
\]
\end{proposition}

\noindent\emph{Proof.} This is the standard bound for an inner product
of length $k$ \citep[Section~3.1]{higham2002}: each product incurs at
most one relative rounding error of size $\mach$, none when $\abs{z_e}$
is a power of two, and any summation order incurs at most $k-1$ more
along the path of each term. Collecting the factors $(1+\delta_j)$ by
Lemma~3.1 of \citet{higham2002} gives the bound. When the data are
differences of potentials, forming each difference adds one rounding
per term, which is the source of $\gamma_{2L}$ in
\Cref{thm:noise-scale}. \hfill$\square$

We span $\ker S$ by the fundamental basis of the reduced row-echelon
form of $S$, computed in exact rational arithmetic and scaled to coprime
integers, so that the basis carries no rounding error of its own.

\subsection*{Data, transform and aggregation}

TECRDB as distributed with component contribution (commit
\texttt{67d3804}) has $4{,}544$ rows; $3{,}912$ report $K'$, $T$ and
pH. Of these, $173$ involve a compound without a structure in the
eQuilibrator cache and were dropped; every remaining reaction balanced
with water and protons. Restricting to grades A--C and $293$--$313$ K
leaves $3{,}276$ measurements from $908$ studies of $418$ distinct
reactions. For a measurement of reaction $e$ at conditions
$c=(\mathrm{pH},I,T,\mathrm{pMg})$,
\[
  d_e \;=\; -RT\ln K'_e \;-\; \sum_i S_{ie}\,\tau_i(c),
\]
where $\tau_i(c)$ is the Legendre transform of compound $i$ computed by
eQuilibrator \citep{beber2022} from its pKa and magnesium-binding
constants; a missing ionic strength is taken as $0.25$ M and a missing
pMg as $14$, as component contribution does \citep{noor2013}. Any
reference-state term in $\tau_i$ is a species-level quantity and
cancels around every cycle (\Cref{lem:species-cancel}), so only the
dependence of $\tau_i$ on each measurement's own conditions reaches a
cycle sum. As a check of the transform, the values recovered at pH~$7$,
$I=0.25$ M and $298.15$ K are $+7.7$ kJ/mol for triose-phosphate
isomerase (DHAP to GAP), $-2.4$ for adenylate kinase (ATP + AMP to 2
ADP), $-3.0$ for enolase (2PG to PEP) and $+3.2$ for glucose-6-phosphate
isomerase (G6P to F6P), in line with tabulated standard values
\citep{alberty1998}.

Measurements repeated within one study are averaged into one value; a
reaction's datum is the mean over its $n_e$ studies, with
$\sigma_e=\sigma/\sqrt{n_e}$ and $\sigma$ the pooled between-study
standard deviation, $\bigl(\sum_e\sum_s(x_{es}-\bar x_e)^2/
\sum_e(n_e-1)\bigr)^{1/2}=2.68$ kJ/mol over the $165$ reactions
measured in at least two studies ($1.49$ kJ/mol for grade A alone).

\subsection*{Calibration on repeated measurements}

For a reaction measured in $n\ge2$ studies, the leave-one-out deviation
$t_s=(x_s-\bar x_{-s})/(\sigma\sqrt{1+1/(n-1)})$ is the standardised
sum of a two-edge cycle whose true sum is zero. Over all grades the
$655$ values exceed the Gaussian quantile $2.576$ in $3.5\%$ of cases
(nominal $1\%$); their $95$th and $99$th percentiles are $2.15$ and
$3.81$, and the largest is $5.26$. For grade A alone the $85$ values
exceed $2.576$ in $1.2\%$ of cases and their $99$th percentile is
$2.58$. Across the $67$ basis cycles of the main analysis the
standardised sums $Z=\hol_z/(\sum_e z_e^2\sigma_e^2)^{1/2}$ have
standard deviation $1.75$, and $10.4\%$ exceed $2.576$.

\begin{table}[h]
\centering\small
\caption{Verdicts on the $67$ basis cycles of the measured-reaction
network (grades A--C) and the $13$ of the grade-A network, by the
quantile used in the data floor, and the number flagged by a fixed
tolerance. No cycle is \textsf{CONSISTENT} under any rule.}
\label{tab:s-real}
\begin{tabular}{@{}llrr@{}}
\toprule
Data & Rule & \textsf{INCONSISTENT} & \textsf{UNDECIDABLE}\\
\midrule
A--C & cycle-local, Gaussian ($2.58$) & $7$ & $60$\\
A--C & cycle-local, replicate-calibrated ($3.81$) & $3$ & $64$\\
A--C & cycle-local, Chebyshev ($10$) & $0$ & $67$\\
A--C & fixed $\eps=10^{-6}$ or $10^{-3}$ kJ/mol & $67$ & ---\\
A--C & fixed $\eps=1$, $5$, $10$ kJ/mol & $54$, $27$, $14$ & ---\\
A    & cycle-local, Gaussian or calibrated ($2.58$) & $2$ & $11$\\
A    & cycle-local, Chebyshev ($10$) & $0$ & $13$\\
\bottomrule
\end{tabular}
\end{table}

\subsection*{The flagged cycles}

The seven cycles flagged under the Gaussian floor in the reference
basis are listed by EC number, with each reaction's coefficient and the
number of studies behind it; ``\&'' marks a single measurement of a
coupled reaction.

\begin{table}[h]
\centering\footnotesize
\caption{Cycles flagged \textsf{INCONSISTENT} (Gaussian data floor,
$\alpha=0.01$) in the reference basis. $\hol$ and $\epsstar$ in kJ/mol.}
\label{tab:s-flagged}
\begin{tabular}{@{}rrrp{0.6\textwidth}@{}}
\toprule
$\hol$ & $\epsstar$ & $Z$ & Reactions: coefficient, EC (studies)\\
\midrule
$-27.77$ & $8.45$ & $-8.46$ & $+1$~2.7.1.90 (1); $-1$~3.1.3.11 (3);
$+1$~3.6.1.1 (6)\\
$-13.77$ & $8.45$ & $-4.20$ & $+1$~2.7.7.4 (3); $-1$~2.7.7.4\&3.6.1.1 (1);
$+1$~3.6.1.1 (6)\\
$23.97$ & $16.18$ & $3.82$ & $+1$~2.1.3.1 (1); $-1$~3.6.1.38 (1);
$-1$~1.1.1.37\&4.1.3.7 (1); $+1$~1.1.1.39 (1); $+1$~4.1.3.7 (1);
$+1$~6.4.1.3 (2)\\
$14.51$ & $11.34$ & $3.30$ & $+1$~1.5.1.15 (1); $-1$~1.4.1.2 (5);
$-1$~1.5.1.5 (1); $+1$~1.4.1.3 (4); $+1$~3.5.4.9 (4)\\
$-14.51$ & $12.28$ & $-3.04$ & $+1$~4.1.2.13 (2); $+1$~3.1.3.1 (2);
$-1$~3.1.3.2 (1); $-1$~2.2.1.2 (1); $-1$~4.1.2.13 (6)\\
$-10.02$ & $9.63$ & $-2.68$ & $+1$~1.1.1.3 (2); $-1$~1.4.1.2 (5);
$+1$~1.4.1.3 (4); $+1$~1.1.1.3 (1)\\
$20.03$ & $19.52$ & $2.64$ & $+1$~2.7.7 (1); $+1$~3.6.1.38 (1);
$+1$~1.1.1.37\&4.1.3.7 (1); $+1$~4.1.1.38 (1); $-1$~1.1.1.39 (1);
$-1$~4.1.3.7 (1); $+1$~3.1.4 (1); $-1$~2.7.9.1 (3); $-2$~3.6.1.1 (6)\\
\bottomrule
\end{tabular}
\end{table}

A repeated EC number within a cycle marks distinct reactions of one
enzyme class (for example, glutamate dehydrogenase with NAD$^+$ and
with NADP$^+$). Over $200$ further fundamental bases, from random
column orders, some cycle was flagged in every basis, the number
flagged ranged from $2$ to $40$, and the intersection of flagged cycles
was always empty. The reactions lying on a flagged cycle in at least
$75\%$ of bases are pyrophosphate--fructose-6-phosphate
1-phosphotransferase (EC~2.7.1.90, $94\%$), methylmalonyl-CoA
carboxyltransferase (2.1.3.1, $88.5\%$), propanoyl-CoA carboxylase
(6.4.1.3, $88.5\%$), methylenetetrahydrofolate dehydrogenase
(NAD$^+$; 1.5.1.15, $80.5\%$) and ADP-glucose phosphodiesterase
(3.1.4, $78\%$); all but one rest on a single study.

\subsection*{The erythrocyte network}

The $38$ genes are HK1, GPI, PFKL, PFKM, ALDOA, TPI1, GAPDH, PGK1, BPGM,
MINPP1, PGAM1, ENO1, PKLR, LDHA, LDHB, G6PD, PGLS, PGD, RPE, RPIA, TKT,
TALDO1, PGM1, AK1, NME1, NME2, GSR, GPX1, GCLC, GSS, ADA, AMPD3, PNP,
APRT, HPRT1, PRPS1, ATP1A1 and ATP2B4. Their UniProtKB/Swiss-Prot
entries annotate $118$ Rhea reactions. We excluded $27$ whose
participants are compound classes without a defined structure (such as
``a D-hexose''), $8$ with generic macromolecular participants, and $9$
with no net chemistry once compartments are merged (pure transport),
leaving $74$, of which $70$ are chemically distinct, over $101$ species.
$S$ has rank $63$, so $\nu=7$, and $19$ reactions lie on a cycle.

One basis cycle is the Rapoport--Luebering shunt: phosphoglycerate
kinase ($3$ studies), phosphoglycerate mutase ($12$),
bisphosphoglycerate mutase ($0$), bisphosphoglycerate phosphatase ($0$)
and ATP hydrolysis by the Ca$^{2+}$ pump ($1$). The other six are cycles
of nucleotide interconversion by adenylate kinase and
nucleoside-diphosphate kinase on (deoxy)nucleotides, with at most one
measured edge each. No cycle has a measured constant on every edge, and
the $31$ erythrocyte reactions with any measurement form a subnetwork
with no cycle at all. Matching to TECRDB is at the level of molecular
connectivity, so an anomer-specific Rhea reaction is matched to
TECRDB's generic one.


\bibliographystyle{apacite}
\bibliography{references}

\end{document}
